Este artigo analisou a evolução do desempenho, das notas de corte e da demanda pelos cursos de Computação da UnB no PAS ao longo de nove subprogramas completos, de 2015--2017 a 2023--2025, sob a perspectiva dos sistemas de concorrência. A partir de uma base extraída e validada dos editais oficiais, com mais de 200 mil registros de candidatos, foi possível acompanhar cada etapa do certame e cada categoria de cota, o que permitiu responder às quatro questões de pesquisa propostas.

Quanto à primeira questão, observou-se que, na UnB como um todo, a distância de desempenho entre o Sistema Universal Estrito e os cotistas diminuiu de forma contínua: o $d$ de Cohen do Argumento Final caiu de 0,69 para 0,21 para Escola Pública e de 0,45 para 0,23 para a cota para Negros, mantendo-se, porém, estatisticamente favorável ao SUE em todos os subprogramas. Essa convergência começou antes da pandemia e da Lei nº 14.723/2023 e foi puxada principalmente pelas categorias PPI. Nos cursos de Computação, a distância partiu de um patamar menor e, no último subprograma, tornou-se estatisticamente indistinguível de zero. Esse resultado deve ser lido com cautela: as categorias de cota na Computação reúnem poucas dezenas de candidatos por ano, o que torna as médias voláteis e faz com que ultrapassagens pontuais entre os grupos sejam esperadas mesmo sem mudança real de desempenho. Assim, os dados indicam uma aproximação entre os sistemas, e não uma inversão consolidada da vantagem do Sistema Universal.

Em relação à segunda questão, as notas de corte revelaram trajetórias distintas entre os cursos. Em Ciência da Computação, o corte do Sistema Universal subiu cerca de 40 pontos de Argumento Final a partir de 2020--2022, acompanhando o crescimento da procura com o mesmo número de vagas, enquanto em Engenharia de Computação ele permaneceu estável. Os cortes das cotas de Escola Pública ficaram próximos da média da coorte, de 100 a 140 pontos abaixo do Sistema Universal, e a cota para Negros ocupou posição intermediária. As chamadas suplementares reduziram a exigência mínima efetiva em todos os grupos dos bacharelados, sobretudo nas categorias com uma única vaga por semestre. Quanto à efetividade do acesso, a reserva de vagas mostrou-se decisiva: nos bacharelados, os cotistas representam entre 21\% e 41\% dos candidatos, mas cerca de metade dos aprovados, próximo da fração de vagas reservadas.

No que se refere à terceira questão, a demanda pelos cursos de Computação cresceu, mas esse crescimento concentrou-se no Sistema Universal. Em Ciência da Computação, a proporção de cotistas entre os candidatos recuou de 38\% para 22\% entre 2017--2019 e 2023--2025. A política de cotas, portanto, preserva a presença dos cotistas entre os aprovados, mas não amplia sua participação na demanda inicial. Na Licenciatura em Computação, a procura foi próxima ou inferior à oferta, de modo que praticamente todos os candidatos foram aprovados e as vagas reservadas não preenchidas migraram para a ampla concorrência; nesse curso, a limitação está na demanda, e não no acesso.

Por fim, quanto à quarta questão, os fatores exógenos mostraram associações específicas, mas não explicam isoladamente as tendências observadas. A queda generalizada das notas de corte em 2018--2020 está ligada principalmente à unificação do ingresso nesse subprograma, e não necessariamente a um pior desempenho no período pandêmico. O aumento da competitividade a partir de 2021 é compatível com a valorização da área de tecnologia, embora os dados não permitam isolar o papel da Inteligência Artificial. A Lei nº 14.723/2023 alterou sobretudo o tipo de vaga ocupada pelos cotistas, sem modificar de forma perceptível o número de cotistas aprovados, e a aproximação de desempenho observada no último subprograma já estava presente nas etapas anteriores à sua aplicação. Com apenas dois subprogramas sob a nova lei, ainda não é possível afirmar mudanças de tendência.

Como trabalhos futuros, propõe-se: (i) acompanhar os próximos subprogramas para avaliar se a aproximação de desempenho na Computação se mantém e quais são os efeitos de médio prazo da Lei nº 14.723/2023; (ii) incorporar dados socioeconômicos dos candidatos, como renda e tipo de escola de origem dos inscritos no Sistema Universal, que não constam dos editais públicos e permitiriam separar efeitos de composição de efeitos de desempenho; (iii) acompanhar a trajetória acadêmica dos ingressantes na graduação, comparando rendimento, retenção e evasão entre cotistas e não cotistas; (iv) incluir Engenharia de Software a partir de dados internos de habilitação da UnB, bem como comparar o PAS com as demais formas de ingresso, como o ENEM e o vestibular; e (v) investigar o efeito da flexibilização de 2022 e da participação em múltiplos subprogramas, o que exigiria identificar o mesmo candidato entre edições do certame.
